\documentclass[11pt,a4paper]{report}

\usepackage[utf8]{inputenc}
\usepackage[T1]{fontenc}
\usepackage[french]{babel}
\usepackage[margin=2.2cm]{geometry}
\usepackage{xurl}
\usepackage{hyperref}
\usepackage{booktabs}
\usepackage{xcolor}
\usepackage{longtable}
\usepackage{enumitem}
\usepackage{fancyhdr}
\usepackage{titlesec}
\usepackage{array}
\usepackage{colortbl}
\usepackage{ragged2e}
\usepackage{graphicx}
\usepackage{tikz}

% ── Palette Copifi (IMMUABLE — issue de la landing en production) ──
\definecolor{copifiobsidian}{HTML}{050505}
\definecolor{copifidark}{HTML}{09090B}
\definecolor{copificardA}{HTML}{121A26}
\definecolor{copificardB}{HTML}{0F1520}
\definecolor{copificardC}{HTML}{0D1825}
\definecolor{copifiemerald}{HTML}{10B981}
\definecolor{copifiemeraldlight}{HTML}{34D399}
\definecolor{copifiemeraldsoft}{HTML}{A7F3D0}
\definecolor{copifisky}{HTML}{0EA5E9}
\definecolor{copifiskylight}{HTML}{38BDF8}
\definecolor{copifiskysoft}{HTML}{BAE6FD}
\definecolor{copifilime}{HTML}{CCFF00}
\definecolor{copifilimehover}{HTML}{B8E600}
\definecolor{copifiblue}{HTML}{3B82F6}
\definecolor{copifizinc400}{HTML}{A1A1AA}
\definecolor{copifizinc500}{HTML}{71717A}
\definecolor{copifitealbrand}{HTML}{11534D}

\hypersetup{
    colorlinks=true,
    linkcolor=copifiemerald,
    urlcolor=copifisky,
    pdftitle={Brief identité visuelle — Copifi},
    pdfauthor={Équipe Copifi},
    pdfsubject={Brief créatif pour designer — SaaS facturation et pilotage}
}

\pagestyle{fancy}
\fancyhf{}
\fancyhead[L]{\small\textit{Brief identité visuelle — Copifi}}
\fancyhead[R]{\small\thepage}
\renewcommand{\headrulewidth}{0.4pt}

\titleformat{\chapter}[display]
{\normalfont\huge\bfseries\color{copifiobsidian}}{\chaptertitlename\ \thechapter}{16pt}{\Huge}

\titleformat{\section}
{\normalfont\Large\bfseries\color{copifiemerald}}
{\thesection}{1em}{}

\newcommand{\hex}[1]{\texttt{\small #1}}
\newcommand{\feat}[1]{\textbf{#1}}
\newcommand{\swatch}[2]{%
    \tikz[baseline=-0.5ex]{\fill[#1] (0,0) rectangle (0.55,0.55);}%
    \hspace{0.35em}#2%
}

\newcommand{\colorrow}[4]{%
    #1 & \hex{#2} & #3 & #4 \\%
}

\begin{document}

% ══════════════════════════════════════════════════════════════════
\begin{titlepage}
    \centering
    \vspace*{1.2cm}
    {\Huge\bfseries\color{copifiobsidian} Brief identité visuelle\par}
    \vspace{0.6cm}
    {\LARGE\color{copifiemerald} Copifi\par}
    \vspace{0.4cm}
    {\large SaaS de facturation conforme 2026 + pilotage d'entreprise\par}
    \vspace{0.3cm}
    {\normalsize Document de cadrage pour designer graphique / identité de marque\par}
    \vspace{1.4cm}
    \rule{0.78\textwidth}{1pt}
    \vspace{1cm}
    \begin{tabular}{>{\bfseries}r p{9.2cm}}
        Client / produit & Copifi (ex FinFlow) \\
        Type de mission & Identité visuelle complète (logo, charte, déclinaisons) \\
        Contrainte majeure & \textbf{Conserver la palette chromatique du site existant} \\
        Référence UI & Landing page \texttt{Welcome.jsx} (juillet 2026) \\
        Cible & TPE, PME, indépendants, auto-entrepreneurs (France) \\
        Version document & 1.0 — Juillet 2026 \\
    \end{tabular}
    \vfill
    \begin{minipage}{0.85\textwidth}
        \centering
        \small\itshape
        Ce document explique le produit, son positionnement, son univers visuel déjà en place
        et les livrables attendus. La palette de couleurs listée au chapitre~\ref{ch:couleurs}
        est \textbf{verrouillée} : le designer crée l'identité \emph{dans} cet univers, sans le remplacer.
    \end{minipage}
\end{titlepage}

\tableofcontents
\cleardoublepage

% ============================================================
\chapter{Mission et périmètre}
% ============================================================

\section{Ce que nous attendons du designer}

Nous souhaitons une \textbf{identité visuelle professionnelle et mémorisable} pour Copifi :
logotype, symbole, déclinaisons, guide d'usage minimal et assets exportés (SVG, PNG, favicon).

Le site web et l'application existent déjà avec une direction visuelle \emph{dark fintech premium}.
\textbf{Votre mission n'est pas de réinventer les couleurs}, mais de :
\begin{itemize}[leftmargin=*]
    \item concevoir un logo et une marque cohérents avec la palette actuelle ;
    \item formaliser une charte graphique (typographies, espacements, usages) ;
    \item produire les déclinaisons nécessaires (web, app, e-mail, print léger) ;
    \item garantir la lisibilité du symbole de 16\,px (favicon) à l'affichage hero.
\end{itemize}

\section{Ce qui est hors périmètre (pour l'instant)}

\begin{itemize}[leftmargin=*]
    \item Refonte complète de l'interface produit (déjà développée en React + Tailwind).
    \item Changement de la palette principale (emerald / sky / lime / obsidian).
    \item Illustrations marketing complexes ou banques d'images custom (hors logo / pictos de base).
\end{itemize}

\section{Livrables attendus}

\begin{enumerate}[leftmargin=*]
    \item \textbf{2 à 3 pistes créatives} avec moodboard court (1 planche par piste).
    \item \textbf{Logo final vectoriel} : symbole + wordmark «\,Copifi\,».
    \item \textbf{Kit de déclinaisons} :
    \begin{itemize}
        \item couleur principale (fond sombre) ;
        \item monochrome blanc (usage principal produit) ;
        \item monochrome noir (fond clair, PDF, e-mails) ;
        \item symbole seul (favicon, icône app, avatar sidebar) ;
        \item wordmark seul si séparable.
    \end{itemize}
    \item \textbf{Favicon} 16$\times$16, 32$\times$32 et app icon 512$\times$512.
    \item \textbf{Mini charte} (PDF 4--8 pages) : couleurs hex, typo recommandée, zone de protection, interdits.
    \item \textbf{Fichiers sources} : Figma ou Adobe Illustrator + exports SVG / PNG @2x.
\end{enumerate}

% ============================================================
\chapter{Le produit Copifi — vue d'ensemble}
% ============================================================

\section{Qu'est-ce que Copifi ?}

Copifi est un \textbf{SaaS B2B français} qui réunit, dans une seule application et un seul abonnement,
deux capacités que les concurrents proposent séparément :

\begin{enumerate}[leftmargin=*]
    \item \feat{Facturation conforme} — devis, factures, paiements, PDF, relances, avec conformité
    à la réforme de la facturation électronique 2026--2027 (Factur-X, Plateforme Agréée DGFiP,
    archivage légal 10 ans).
    \item \feat{Pilotage d'entreprise} — tableau de bord dirigeant : trésorerie, marge, CAC, LTV, DSO,
    alertes automatiques, analyses IA mensuelles et copilote financier contextuel.
\end{enumerate}

\section{Promesse et baseline}

\begin{table}[h]
\centering
\caption{Messages clés de marque (déjà utilisés sur le site)}
\begin{tabular}{@{}p{3.2cm}p{10.5cm}@{}}
\toprule
\textbf{Élément} & \textbf{Contenu} \\
\midrule
Baseline hero & \emph{Pilotez, facturez, décidez.} \\
Sous-titre & Bien plus qu'un logiciel de facturation. \\
Accroche différenciation & Les logiciels de facturation s'arrêtent à la facture. Nous, on va jusqu'à la décision. \\
Promesse conformité & La loi change. Pas vos habitudes. \\
Promesse pilotage & Vos chiffres ne dorment plus dans un tableur. \\
CTA principal & Démarrer l'essai gratuit (14 jours, sans carte bancaire). \\
\bottomrule
\end{tabular}
\end{table}

\section{Public cible}

\begin{itemize}[leftmargin=*]
    \item \textbf{Indépendants et freelances} — besoin de facturer vite et d'encaisser sans friction.
    \item \textbf{TPE / PME} — 1 à 20 salariés, dirigeant seul face aux chiffres, sans DAF dédié.
    \item \textbf{Auto-entrepreneurs assujettis à la TVA} — concernés par la réforme 2026.
    \item \textbf{Profil psychologique} : entrepreneur pragmatique, sensible à la conformité légale,
    allergique aux tableurs Excel, en quête de clarté sans jargon comptable.
\end{itemize}

\section{Concurrents implicites et différenciation}

\begin{table}[h]
\centering
\caption{Positionnement}
\begin{tabular}{@{}p{4.5cm}p{9.2cm}@{}}
\toprule
\textbf{Catégorie} & \textbf{Limite} \\
\midrule
Logiciels de facturation classiques & S'arrêtent à l'émission de facture ; pas de pilotage financier. \\
Tableurs Excel / Google Sheets & Manuels, erreurs, pas de conformité PA / Factur-X. \\
Outils BI génériques & Complexes, pas intégrés au cycle commercial. \\
\textbf{Copifi} & \textbf{Facturation + pilotage + IA, un seul abonnement, données auto-alimentées par chaque facture.} \\
\bottomrule
\end{tabular}
\end{table}

\section{Les deux piliers produit (système visuel à refléter)}

L'identité doit permettre de \textbf{distinguer visuellement} les deux piliers, comme sur le site :

\begin{table}[h]
\centering
\caption{Sémantique des deux piliers}
\rowcolors{2}{white}{gray!6}
\begin{tabular}{@{}p{2.8cm}p{5.5cm}p{5.5cm}@{}}
\toprule
 & \textbf{Pilier 1 — Facturation} & \textbf{Pilier 2 — Pilotage} \\
\midrule
Couleur UI & Emerald (vert) & Sky (bleu ciel) \\
Thème & Conformité, légal, confiance & Data, KPIs, intelligence \\
Fonctions clés & Devis, factures, Factur-X, PA, archivage & Trésorerie, marge, CAC, LTV, alertes, IA \\
Message & Facturer, sans faute & Piloter, sans tableur \\
\bottomrule
\end{tabular}
\end{table}

% ============================================================
\chapter{Fonctionnalités détaillées du SaaS}
% ============================================================

\section{Module facturation (Pilier 1)}

\begin{itemize}[leftmargin=*]
    \item \feat{Gestion clients (tiers)} — fiches clients, coordonnées, pays, numéro d'immatriculation.
    \item \feat{Catalogue produits / services} — articles avec prix, descriptions, images, activation.
    \item \feat{Devis} — création, envoi par e-mail (PDF), suivi d'ouverture (pixel tracking),
    statuts (brouillon, envoyé, accepté, refusé, expiré), conversion en facture en un clic.
    \item \feat{Factures} — numérotation légale automatique, lignes avec remises, escompte financier,
    multi-devises, prestations, catégories d'opération, export PDF.
    \item \feat{Conformité électronique} — génération Factur-X (PDF + XML), transmission Plateforme Agréée,
    statuts CDAR, archivage inaltérable 10 ans, piste d'audit.
    \item \feat{Paiements} — suivi des encaissements, relances, tableau de bord recouvrement,
    escompte sur paiement anticipé.
    \item \feat{Factures fournisseurs} — saisie des achats (module en cours d'enrichissement).
    \item \feat{Paramètres entreprise} — logo, coordonnées, TVA, mandat de facturation électronique,
    destinations de livraison.
\end{itemize}

\section{Module pilotage financier (Pilier 2)}

\begin{itemize}[leftmargin=*]
    \item \feat{Dashboard CFO} — KPI en temps réel : CA, charges, marge, trésorerie.
    \item \feat{Indicateurs avancés} — CAC (coût d'acquisition client), LTV (valeur vie client),
    DSO (délai de paiement), ratios de rentabilité.
    \item \feat{Saisie mensuelle} — complément de données pour analyse (charges fixes, etc.).
    \item \feat{Alertes automatiques} — marge négative, charges $>$ 70\,\% du CA, retards d'encaissement.
    \item \feat{Analyse IA mensuelle} — synthèse en français clair de la santé de l'activité.
    \item \feat{Copilote IA} — assistant conversationnel connaissant les chiffres de l'entreprise
    (Groq / Llama 3.3 70B).
    \item \feat{Graphiques} — courbes CA vs charges, donuts, cartes KPI animées (Recharts).
\end{itemize}

\section{Acquisition, monétisation et parcours utilisateur}

\begin{itemize}[leftmargin=*]
    \item \feat{Landing page marketing} — page d'accueil longue, sections éducatives, CTA essai gratuit.
    \item \feat{Assistant IA landing} — chat commercial pour répondre aux questions (réforme, tarifs, produit).
    \item \feat{Authentification} — inscription, connexion, Google OAuth, vérification e-mail.
    \item \feat{Abonnement Stripe} — plans Indépendant (15\,€/mois HT) et Pro (49\,€/mois HT),
    essai 14 jours sans carte, checkout et webhooks.
    \item \feat{Espace admin} — gestion utilisateurs, MRR, suspension de comptes.
\end{itemize}

\section{Contexte réglementaire (argument marketing central)}

\begin{table}[h]
\centering
\caption{Calendrier réforme facturation électronique (France)}
\begin{tabular}{@{}llp{8.5cm}@{}}
\toprule
\textbf{Date} & \textbf{Obligation} & \textbf{Impact communication} \\
\midrule
1\textsuperscript{er} sept.\ 2026 & Réception obligatoire & Toutes les entreprises assujetties TVA doivent pouvoir recevoir des factures électroniques. \\
1\textsuperscript{er} sept.\ 2027 & Émission obligatoire (TPE/PME) & Les petites structures devront émettre via une Plateforme Agréée. \\
\bottomrule
\end{tabular}
\end{table}

Copifi transforme cette contrainte en \textbf{opportunité marketing} : la conformité est gérée
automatiquement, l'utilisateur continue de facturer comme avant.

% ============================================================
\chapter{Univers visuel existant — direction artistique}
\label{ch:univers}
% ============================================================

\section{Style global : \emph{Dark Fintech Premium}}

L'interface actuelle suit une esthétique inspirée des sites primés (Awwwards) :
\begin{itemize}[leftmargin=*]
    \item fond \textbf{obsidian} profond, quasi noir ;
    \item lueurs néon \textbf{diffuses} (emerald et sky), jamais agressives ;
    \item cartes en \textbf{glassmorphism} : bordures fines blanches semi-transparentes, fonds dégradés sombres ;
    \item coins très arrondis (\texttt{border-radius} 2\,rem typique) ;
    \item animations fluides avec courbe \texttt{cubic-bezier(0.22, 1, 0.36, 1)} (effet ressort amorti) ;
    \item typographie éditoriale : titres massifs + mots-clés en \textbf{serif italique} avec dégradé.
\end{itemize}

\section{Registre émotionnel à transmettre}

\begin{table}[h]
\centering
\begin{tabular}{@{}p{4cm}p{5cm}p{5cm}@{}}
\toprule
\textbf{Oui} & \textbf{Non} \\
\midrule
Clarté, sérénité, modernité & Froideur bancaire institutionnelle \\
Confiance, sérieux, innovation maîtrisée & Startup gadget / crypto flashy \\
Premium accessible, entrepreneur-friendly & Corporate costume-cravate \\
Data-driven, précision & Illustrations stock cliché finance \\
\bottomrule
\end{tabular}
\end{table}

\section{Références d'ambiance (non copies)}

\begin{itemize}[leftmargin=*]
    \item SaaS fintech dark mode (Linear, Stripe docs dark, Vercel marketing).
    \item Sites Awwwards : micro-interactions, profondeur, glass cards.
    \item Dashboards financiers épurés : lisibilité des chiffres, hiérarchie forte.
\end{itemize}

% ============================================================
\chapter{Palette chromatique — CONSERVER TELLE QUELLE}
\label{ch:couleurs}
% ============================================================

\section{Règle impérative}

\textbf{Les couleurs ci-dessous sont celles du site en production.}
Le designer \textbf{ne doit pas les remplacer} par une autre palette.
Il peut proposer des nuances \emph{à l'intérieur} de cette famille (opacités, dégradés),
mais les teintes principales restent fixes.

\section{Fonds et surfaces}

\begin{table}[h]
\centering
\caption{Couleurs de fond}
\rowcolors{2}{white}{gray!5}
\begin{tabular}{@{}clll@{}}
\toprule
\textbf{Échantillon} & \textbf{Hex} & \textbf{Nom} & \textbf{Usage} \\
\midrule
\colorrow{\swatch{copifiobsidian}{}}{050505}{Obsidian}{Fond principal page, body}
\colorrow{\swatch{copifidark}{}}{09090B}{Obsidian alt.}{Zones app (Tailwind \texttt{obsidian})}
\colorrow{\swatch{copificardA}{}}{121A26}{Card A}{Dégradé cartes (haut)}
\colorrow{\swatch{copificardB}{}}{0F1520}{Card B}{Dégradé cartes (milieu)}
\colorrow{\swatch{copificardC}{}}{0D1825}{Card C}{Dégradé cartes (bas)}
\bottomrule
\end{tabular}
\end{table}

\noindent Surfaces glass : blanc à \textbf{2--8\,\%} d'opacité (\texttt{rgba(255,255,255,0.02--0.08)}).
Bordures glass : blanc à \textbf{10\,\%} (\texttt{border-white/10}).

\section{Accent Pilier 1 — Emerald (facturation / conformité)}

\begin{table}[h]
\centering
\caption{Famille emerald — NE PAS REMPLACER}
\rowcolors{2}{white}{gray!5}
\begin{tabular}{@{}clll@{}}
\toprule
\textbf{Échantillon} & \textbf{Hex} & \textbf{Tailwind} & \textbf{Usage} \\
\midrule
\colorrow{\swatch{copifiemerald}{}}{10B981}{emerald-500}{Badges, icônes, sélection texte, CTA secondaires}
\colorrow{\swatch{copifiemeraldlight}{}}{34D399}{emerald-400}{Contours logo, hover cartes, labels}
\colorrow{\swatch{copifiemeraldsoft}{}}{A7F3D0}{emerald-200}{Dégradés titres (avec sky-300)}
\bottomrule
\end{tabular}
\end{table}

Lueurs d'ambiance : \texttt{rgba(16,185,129,0.15--0.25)} en radial blur.

\section{Accent Pilier 2 — Sky (pilotage / data)}

\begin{table}[h]
\centering
\caption{Famille sky — NE PAS REMPLACER}
\rowcolors{2}{white}{gray!5}
\begin{tabular}{@{}clll@{}}
\toprule
\textbf{Échantillon} & \textbf{Hex} & \textbf{Tailwind} & \textbf{Usage} \\
\midrule
\colorrow{\swatch{copifisky}{}}{0EA5E9}{sky-500}{Section pilotage, KPIs, graphiques}
\colorrow{\swatch{copifiskylight}{}}{38BDF8}{sky-400}{Courbes graphiques, hovers}
\colorrow{\swatch{copifiskysoft}{}}{BAE6FD}{sky-300}{Dégradés titres, accents doux}
\bottomrule
\end{tabular}
\end{table}

Lueurs : \texttt{rgba(56,189,248,0.1--0.15)}.

\section{CTA et highlight}

\begin{table}[h]
\centering
\caption{Couleur d'action primaire}
\rowcolors{2}{white}{gray!5}
\begin{tabular}{@{}clll@{}}
\toprule
\textbf{Échantillon} & \textbf{Hex} & \textbf{Nom} & \textbf{Usage} \\
\midrule
\colorrow{\swatch{copifilime}{}}{CCFF00}{Lime CTA}{Boutons «\,Démarrer l'essai gratuit\,»}
\colorrow{\swatch{copifilimehover}{}}{B8E600}{Lime hover}{État survol CTA}
\bottomrule
\end{tabular}
\end{table}

Cette lime est \textbf{signature conversion} : haute visibilité sur fond sombre, texte noir dessus.

\section{Accents complémentaires}

\begin{itemize}[leftmargin=*]
    \item \textbf{Bleu} \hex{3B82F6} — courbes secondaires graphiques, dégradés internes cartes.
    \item \textbf{Teal marque historique} \hex{11534D} — présent dans la sidebar app (dégradé radial) ;
    peut subsister en accent secondaire mais \textbf{ne remplace pas} emerald/sky de la landing.
    \item \textbf{Ambre} \texttt{\#F59E0B} — alertes «\,en attente\,», warnings légers.
    \item \textbf{Rose} — messages d'erreur flash uniquement.
\end{itemize}

\section{Texte et neutres}

\begin{table}[h]
\centering
\rowcolors{2}{white}{gray!5}
\begin{tabular}{@{}clll@{}}
\toprule
\textbf{Échantillon} & \textbf{Hex} & \textbf{Rôle} & \textbf{Usage} \\
\midrule
\colorrow{\swatch{white}{}}{FFFFFF}{Titres}{H1, H2, H3, navigation}
\colorrow{\swatch{copifizinc400}{}}{A1A1AA}{Corps}{Paragraphes, descriptions (zinc-400)}
\colorrow{\swatch{copifizinc500}{}}{71717A}{Labels}{Kickers, métadonnées (zinc-500/600)}
\bottomrule
\end{tabular}
\end{table}

\section{Dégradé signature des titres éditoriaux}

Mots en italique serif sur la landing :
\[
\texttt{from-emerald-200}\ (#A7F3D0) \;\rightarrow\; \texttt{to-sky-300}\ (#BAE6FD)
\]
Appliqué en \texttt{background-clip: text} sur fond sombre.

% ============================================================
\chapter{Typographie}
% ============================================================

\section{Polices actuellement utilisées}

\begin{table}[h]
\centering
\caption{Hiérarchie typographique}
\begin{tabular}{@{}p{3.5cm}p{4cm}p{7cm}@{}}
\toprule
\textbf{Rôle} & \textbf{Police} & \textbf{Usage} \\
\midrule
UI / corps & Inter & Application, formulaires, tableaux, boutons \\
Titres app & Syne (+ Inter fallback) & En-têtes modules facturation (\texttt{font-display}) \\
Éditorial landing & Instrument Serif & Accroches italiques : \emph{décidez.}, \emph{sans friction.} \\
Labels techniques & IBM Plex Mono & Badges, stats, compte à rebours, marquee trust \\
\bottomrule
\end{tabular}
\end{table}

\section{Recommandations pour le wordmark «\,Copifi\,»}

\begin{itemize}[leftmargin=*]
    \item Lisible, moderne, légèrement tech sans agressivité futuriste.
    \item Cohérent avec Instrument Serif (landing) \textbf{ou} Inter/Syne (app) — le designer choisit
    une personnalité propre au logo tout en restant compatible.
    \item Le «\,C\,» peut porter une légère personnalité (référence au symbole) sans devenir illisible.
    \item Éviter les serifs classiques type Times : trop institutionnel banque.
\end{itemize}

\section{Principes de mise en page}

\begin{itemize}[leftmargin=*]
    \item Titres landing : \texttt{tracking-tighter}, tailles XXL (48--72\,px desktop).
    \item Kickers : uppercase, letterspacing 0.18--0.28\,em, 11\,px, mono.
    \item Chiffres KPI : chasse tabulaire (\texttt{tabular-nums}), police mono.
\end{itemize}

% ============================================================
\chapter{Logo et symbole — pistes et contraintes}
% ============================================================

\section{État actuel (provisoire)}

Un logo temporaire existe : icône \emph{document à coin plié} + courbe de flux (emerald $\rightarrow$ lime),
wordmark «\,Copifi\,» en Instrument Serif, avec animation au survol (micro-rotations, glow emerald).
\textbf{Ce logo est un placeholder technique}, pas la identité finale.

\section{Concepts à explorer (liberté créative)}

Le designer est libre d'interpréter, en restant dans l'univers décrit :

\begin{itemize}[leftmargin=*]
    \item \textbf{Flux / courbe} — croissance maîtrisée, pas une flèche stock market cliché.
    \item \textbf{Document / facture stylisée} — lien avec la facturation, forme épurée.
    \item \textbf{Pilotage / jauge / signal} — clarté, mesure, décision.
    \item \textbf{Monogramme C} — si lisible en 16\,px.
    \item \textbf{Combinaison symbole + wordmark} — horizontal (nav) et empilé (carré app).
\end{itemize}

\section{À éviter absolument}

\begin{itemize}[leftmargin=*]
    \item Calculatrice, pièces de monnaie, symboles \$ € £.
    \item Graphiques en barres ou camemberts littéraux (trop BI générique).
    \item Bouclier ou cadenas seul (cybersécurité, pas clarté financière).
    \item Costume-cravate, poignée de main (corporate banque).
    \item Néons agressifs type crypto — le glow Copifi reste \textbf{subtil et diffus}.
    \item Logo surchargé illisible en favicon 16\,px.
    \item Palette hors charte (violet, orange vif, rose néon comme couleur primaire).
\end{itemize}

\section{Contraintes techniques logo}

\begin{itemize}[leftmargin=*]
    \item 100\,\% vectoriel (SVG source obligatoire).
    \item Lisible sur \textbf{fond \hex{050505}} (usage principal) et fond blanc (e-mails, PDF).
    \item Version monochrome blanche pour sidebar et footer.
    \item Zone de protection = hauteur du symbole minimum de chaque côté.
    \item Taille minimale web : 24\,px hauteur symbole ; favicon 16\,px reconnaissable.
\end{itemize}

% ============================================================
\chapter{Composants UI — langage graphique à respecter}
% ============================================================

\section{Vocabulaire visuel des composants}

\begin{description}[leftmargin=3.5cm, style=nextline]
    \item[GlassCard] Carte avec bordure \texttt{white/10}, fond dégradé \hex{121A26}$\rightarrow$\hex{0D1825},
    ombre profonde, hover avec halo emerald ou sky.
    \item[Nav pill] Barre flottante centrée, \texttt{rounded-full}, glass + blur.
    \item[Badge] Pill coloré semi-transparent (emerald ou sky), texte mono uppercase.
    \item[CTA primaire] Pill lime \hex{CCFF00}, texte noir, glow néon au hover.
    \item[CTA secondaire] Bordure glass, fond \texttt{white/5--10\%}, texte blanc.
    \item[Marquee] Bandeau défilant trust items, mono, zinc-500, séparateurs emerald.
\end{description}

\section{Iconographie}

\begin{itemize}[leftmargin=*]
    \item Style : \textbf{SVG outline}, trait 1.8--2\,px, coins arrondis.
    \item Pas d'illustrations 3D ni de pictos remplis lourds.
    \item Couleur héritée du contexte (emerald = facturation, sky = pilotage).
    \item Icônes clés : document, bouclier, jauge, cloche, cerveau (IA), flèche, sparkles.
\end{itemize}

\section{Motion design (référence, pas à produire en logo statique)}

\begin{itemize}[leftmargin=*]
    \item Easing signature : \texttt{cubic-bezier(0.22, 1, 0.36, 1)}.
    \item Entrées scroll : fade + translate Y, délais échelonnés.
    \item Hover cartes : élévation $-4$\,px + halo coloré.
    \item Logo nav : micro-animation au survol (référence comportementale).
\end{itemize}

% ============================================================
\chapter{Contexts d'application de la marque}
% ============================================================

\begin{table}[h]
\centering
\caption{Touchpoints où le logo et la charte seront utilisés}
\begin{tabular}{@{}p{4cm}p{4cm}p{6cm}@{}}
\toprule
\textbf{Support} & \textbf{Fond} & \textbf{Contrainte} \\
\midrule
Landing page (nav) & Obsidian \hex{050505} & Logo horizontal, animation hover \\
Sidebar application & Dégradé teal sombre & Symbole 40$\times$40\,px + label \\
Favicon navigateur & Obsidian ou transparent & 16$\times$16 reconnaissable \\
Page connexion & Split : photo + glass clair & Logo couleur sur fond clair \\
E-mails transactionnels & Blanc ou gris clair & Version noir ou couleur \\
Factures PDF & Blanc & Monochrome ou couleur sobre \\
Stripe Checkout & Blanc (Stripe) & Logo compact \\
Réseaux sociaux & Obsidian & Carré 1080$\times$1080 \\
\bottomrule
\end{tabular}
\end{table}

% ============================================================
\chapter{Structure de la landing page (référence contenu)}
% ============================================================

Pour que le designer comprenne le \textbf{volume et le rythme visuel} du site :

\begin{enumerate}[leftmargin=*]
    \item \textbf{Hero} — titre XXL, dashboard 3D flottant, 2 cartes animées (facture / marge).
    \item \textbf{Marquee} — 12 garanties défilantes.
    \item \textbf{La différence} — 2 piliers côte à côte (emerald vs sky).
    \item \textbf{Stats} — 4 chiffres animés.
    \item \textbf{Facturation} — cycle devis $\rightarrow$ facture $\rightarrow$ paiement.
    \item \textbf{Conformité 2026} — compte à rebours, timeline 2026/2027, 3 cartes Factur-X / PA / archivage.
    \item \textbf{Pilotage} — checklist + graphique + 3 cartes (KPIs, alertes, copilote IA).
    \item \textbf{Assistant IA} — chat widget + CTA.
    \item \textbf{Tarifs} — 2 plans (15\,€ / 49\,€).
    \item \textbf{FAQ} — 6 questions accordion.
    \item \textbf{CTA final} — urgence réforme + essai gratuit.
    \item \textbf{Footer} — logo, liens légaux.
\end{enumerate}

% ============================================================
\chapter{Tarification (cohérence communication)}
% ============================================================

\begin{table}[h]
\centering
\caption{Offres commerciales}
\begin{tabular}{@{}llp{7.5cm}@{}}
\toprule
\textbf{Plan} & \textbf{Prix} & \textbf{Contenu} \\
\midrule
Indépendant & 15\,€/mois HT & 30 factures/mois, devis illimités, conformité 2026, archivage, dashboard CA/marge/trésorerie. \\
Pro (populaire) & 49\,€/mois HT & Factures illimitées, pilotage avancé (CAC, LTV, DSO), alertes + copilote IA, multi-devises, support prioritaire. \\
\bottomrule
\end{tabular}
\end{table}

\noindent\textbf{Règle marketing :} les deux piliers sont inclus dans \emph{chaque} plan — ce n'est pas une option.

% ============================================================
\chapter{North star — synthèse pour le designer}
% ============================================================

\begin{quote}
\large\itshape
Copifi, c'est la clarté financière enfin accessible : facturer conforme à la loi,
piloter son entreprise comme un dirigeant équipé, décider sans tableur ni jargon.
L'identité visuelle doit dire \textbf{confiance + modernité + intelligence},
dans un univers \textbf{dark premium} aux accents \textbf{emerald}, \textbf{sky} et \textbf{lime},
sans clichés comptables.
\end{quote}

\vspace{1em}
\noindent\textbf{Checklist avant livraison :}
\begin{itemize}[leftmargin=*]
    \item[\checkmark] Palette \hex{050505} / emerald / sky / \hex{CCFF00} respectée ?
    \item[\checkmark] Logo lisible en 16\,px ?
    \item[\checkmark] Versions fond sombre + fond clair ?
    \item[\checkmark] Cohérent avec les deux piliers (facturation vs pilotage) ?
    \item[\checkmark] Évite les clichés listés au chapitre~8 ?
    \item[\checkmark] Fichiers vectoriels + mini charte fournis ?
\end{itemize}

\vfill
\begin{center}
\rule{0.5\textwidth}{0.4pt}\\[0.5em]
\small
Contact projet : équipe Copifi — Document v1.0 — Juillet 2026\\
Compilable avec \texttt{pdflatex} ou \texttt{xelatex}
\end{center}

\end{document}
